\documentclass[11pt]{article}
\usepackage[margin=1in]{geometry}
\usepackage{amsmath,amssymb,amsthm,mathtools}
\usepackage{graphicx}
\usepackage[hidelinks]{hyperref}

\newtheorem{theorem}{Theorem}
\newtheorem{remark}[theorem]{Remark}
\newcommand{\T}{(0,2\pi)}
\newcommand{\R}{\mathbb R}
\setlength{\emergencystretch}{2em}

\title{No universal weighted Sobolev--SVD tail bound\\
\large A smooth rank-two Fourier counterexample}
\author{Autonomous mathematical discovery run \texttt{fa\_banach\_001}}
\date{17 August 2026}

\begin{document}
\maketitle

\noindent\textbf{Status.}
Candidate full negative answer, likely valid pending expert review.

\section{Source question}

Ali and Nouy compare the ordinary $L^2$ singular value decomposition
\[
 u=\sum_{k\geq1}\sigma_k^{00}\psi_k\otimes\phi_k
\]
with the SVD obtained by viewing the same kernel as a Hilbert--Schmidt
operator $L^2(\Omega_2)\to H^1(\Omega_1)$; its singular values are denoted
$\sigma_k^{10}$.  Immediately after Proposition 2.1 of \cite{AN}, on PDF
page 3, they ask whether one can bound
\begin{equation}\label{eq:question}
 \sum_{k>r}(\sigma_k^{00})^2\|\psi_k\|_{H^1(\Omega_1)}^2
 \lesssim
 \sum_{k>r}(\sigma_k^{10})^2\gamma(k)
\end{equation}
for some sequence $\gamma(k)$.

\begin{center}
\includegraphics[width=0.96\textwidth]{figures/open_problem_crop.png}
\end{center}

We use the only nontrivial uniform interpretation: the finite-valued sequence
$\gamma$ and the implicit constant are fixed independently of $u$.  If either
may depend arbitrarily on $u$, the assertion carries no uniform content.

\section{Proof intuition}

Mix one very high Fourier mode with the constant mode by a $45$-degree
rotation.  The resulting two functions remain orthogonal in $L^2$, so they
can be prescribed as the unique left singular functions of an $L^2$-SVD.
In $H^1$, however, almost all of their growing derivative energy points in
the same direction.

The strong-norm SVD detects that common direction and places the growing
energy in its first singular value.  Its second singular value stays bounded.
After truncating the first term, the source's weighted $L^2$ tail still
contains half of the high-frequency derivative energy, while the proposed
right side contains only the bounded second strong-norm singular value.  The
failure already occurs at index two, so no growth of later weights can help.

\section{Counterexample}

Equip $H^1(\T)$ with
\[
 \langle f,g\rangle_{H^1}=\int_0^{2\pi}(fg+f'g')\,dx.
\]
Let
\[
 e_0(x)=(2\pi)^{-1/2},
 \qquad e_N(x)=\pi^{-1/2}\cos(Nx),
\]
so that $e_0,e_N$ are $L^2$-orthonormal,
$\|e_0'\|_2=0$, and $\|e_N'\|_2=N$.  Put
\[
 \psi_1^{(N)}=\frac{e_N+e_0}{\sqrt2},
 \qquad
 \psi_2^{(N)}=\frac{e_N-e_0}{\sqrt2}.
\]
Choose any smooth $L^2(\T)$-orthonormal functions $\phi_1,\phi_2$.

\begin{theorem}\label{thm:main}
There is no finite-valued sequence $\gamma:\mathbb N\to[0,\infty)$ and no
constant $C<\infty$ for which \eqref{eq:question}, with $\lesssim$ replaced by
$\leq C$, holds for every $u\in H^1(\T^2)$ whose $L^2$-SVD converges in
$H^1(\T^2)$ and every $r\geq0$.  The failure already occurs for smooth
rank-two functions and $r=1$.
\end{theorem}

\begin{proof}
For $N\geq1$, define
\begin{equation}\label{eq:uN}
 u_N(x,y)=2\psi_1^{(N)}(x)\phi_1(y)
          +\psi_2^{(N)}(x)\phi_2(y).
\end{equation}
This is smooth and has finite rank.  Since the two $\psi_j^{(N)}$ and the two
$\phi_j$ are $L^2$-orthonormal and the coefficients $2>1$ are distinct,
\eqref{eq:uN} is its unique $L^2$-SVD up to signs.  Hence
\[
 \sigma_1^{00}=2,\qquad \sigma_2^{00}=1,
\]
and, at $r=1$, the left side of \eqref{eq:question} is
\begin{equation}\label{eq:left}
 (\sigma_2^{00})^2\|\psi_2^{(N)}\|_{H^1}^2
 =1+\frac{N^2}{2}.
\end{equation}

It remains to compute the $H^{(1,0)}$ singular values.  The $H^1$ Gram
matrix of the two left $L^2$ singular functions is
\begin{equation}\label{eq:gram}
 G_N=
 \begin{pmatrix}
  1+N^2/2&N^2/2\\
  N^2/2&1+N^2/2
 \end{pmatrix}.
\end{equation}
Indeed, the $L^2$ contribution is the identity, while both derivatives equal
$e_N'/\sqrt2$.

As an operator $L^2(\T)\to H^1(\T)$, $u_N$ sends $\phi_1$ to
$2\psi_1^{(N)}$, sends $\phi_2$ to $\psi_2^{(N)}$, and annihilates their
orthogonal complement.  Therefore the two nonzero squared
$H^{(1,0)}$ singular values are the eigenvalues
$\lambda_+(N)\geq\lambda_-(N)>0$ of
\begin{equation}\label{eq:AN}
 A_N=\operatorname{diag}(2,1)G_N\operatorname{diag}(2,1)
 =\begin{pmatrix}
   4+2N^2&N^2\\
   N^2&1+N^2/2
  \end{pmatrix}.
\end{equation}
Now
\[
 \det A_N=4(1+N^2),
 \qquad
 \lambda_+(N)\geq (A_N)_{11}=4+2N^2.
\]
It follows that
\begin{equation}\label{eq:small}
 (\sigma_2^{10})^2=\lambda_-(N)
 =\frac{\det A_N}{\lambda_+(N)}
 \leq\frac{4(1+N^2)}{4+2N^2}<2.
\end{equation}
There are no further nonzero singular values.  Thus the right side of the
purported bound at $r=1$ is at most
\[
 C\gamma(2)(\sigma_2^{10})^2<2C\gamma(2),
\]
which is independent of $N$.  This contradicts \eqref{eq:left} as
$N\to\infty$.
\end{proof}

\begin{remark}[Symmetric companion]
Interchanging $x$ and $y$ gives the same counterexample to the analogous
bound involving the $H^{(0,1)}$ singular values $\sigma_k^{01}$ and the
$H^1$ norms of the right $L^2$ singular functions.
\end{remark}

\section{Verification, limitations, and novelty}

The accompanying script verifies the Gram matrix, determinant, trace, the
uniform bound $\lambda_-<2$, and the diverging ratio for several exact values
of $N$.  These checks are not part of the proof.

The theorem settles the displayed question negatively under its natural
uniform reading, even for smooth finite-rank functions on a fixed domain.  It
does not exclude estimates when $\gamma$ or the comparison constant depends
on the particular function, nor estimates under additional assumptions such
as spectral localization or $H^1$-incoherence.  Those nonuniform or restricted
versions are consistent with the source's caveat “without further
assumptions.”

A bounded novelty search on 17 August 2026 covered the four run indexes, the
local source, the current arXiv and journal records, the exact displayed
question, title/arXiv-id queries, and Sobolev--SVD weighted-tail variants.  No
later answer or this rank-two Fourier family was found.  Novelty confidence is
moderate pending specialist review.

\paragraph{Human review recommendation.}
Check the normalization of $e_N$, the Gram matrix \eqref{eq:gram}, the
two-by-two strong-norm singular-value reduction, and the uniform
interpretation of $\gamma$ and the implicit constant.

\enlargethispage{2.5\baselineskip}
\begin{thebibliography}{9}
\small
\raggedright
\bibitem{AN}
M. Ali and A. Nouy,
\emph{Singular Value Decomposition in Sobolev Spaces: Part II},
Z. Anal. Anwend. 40 (2021), 187--216; arXiv:1912.11293.
\end{thebibliography}

\end{document}
